% !TeX program = lualatex
\documentclass[a4paper,11pt]{article}
\usepackage{luatexja-fontspec}
\usepackage{amsmath,amssymb,amsthm}
\usepackage[unicode]{hyperref}

\newcommand{\dfix}{\bigcirc}
\setlength{\parindent}{0pt}
\setlength{\parskip}{1\baselineskip}

\title{二値截断圏と中道不動点：冗（Amphigory）の数学的準備}
\author{千葉将臣}
\date{2026年7月7日}

\begin{document}

\maketitle

\begin{abstract}
本稿は、冗（Amphigory）の本格的定式化に先立ち、その基礎概念である二値截断圏 $\mathcal{B}_{\mathrm{val}}$ と中道不動点 $\dfix$ を、既存の論理学・圏論・不動点理論との接続において整理する準備論文である。ここでは、二値截断圏を古典論理の二値性がより広い論理空間へ埋め込まれる部分圏として解釈し、中道不動点を $E_\infty$ における最大固定点として説明する。これにより、冗において導入される視点数、継続残余、無記、および空への収束が、一般的な数学知識との差分としてどの位置に現れるかを明確化する。
\end{abstract}

\section{はじめに}
冗は、Freyd と Scedrov による古典的な Allegory（寓）の公理系に、視点数（PoV: Point of View）および継続残余 $\rho$ を導入することで、静的な関係代数を動的な証明探索・評価過程へ拡張しようとする枠組みである。

その本格的な定式化では、古典的な「成立するか・しないか」という二値的評価を、二値截断圏 $\mathcal{B}_{\mathrm{val}}$ として取り出す。また、非停止プロセスや自己言及的な無限退行は、単なるエラーとして排除されるのではなく、残余空間へ隔離され、極限において中道不動点 $\dfix$ へ収束するものとして扱われる。

本稿の目的は、この二つの基礎概念、すなわち「二値截断圏」と「中道不動点 $\dfix$」を、一般的な数学知識の枠内で説明することである。これにより、冗の新規性が既存数学からの断絶ではなく、既存の論理学・圏論・不動点理論に対する差分として位置づけられることを示す。

\section{二値截断圏：古典論理の直観主義論理への埋め込み}

\subsection{定義}
冗の定義において、二値截断圏 $\mathcal{B}_{\mathrm{val}}$ は、「射が『成立するか・しないか』という二値評価へ截断される部分圏」として導入される。
すなわち、対象間の射が真理値 $\{\top, \bot\}$ によって評価される部分圏である。

ここで重要なのは、$\mathcal{B}_{\mathrm{val}}$ が冗全体ではなく、そのうち古典的な二値評価だけを取り出した截断であるという点である。冗では、この外側に視点数、継続残余、無記の状態が存在する。

\subsection{古典論理の直観主義論理への埋め込みとしての解釈}
古典論理（二値論理）は、直観主義論理へ埋め込むことができる。代表的な例として、Gödel--Gentzen の二重否定翻訳が知られている。ここでは、古典命題 $P$ を直観主義論理における $\neg\neg P$ に対応させる。

二値截断圏 $\mathcal{B}_{\mathrm{val}}$ は、この埋め込み構造を圏論的に表現したものと解釈できる。すなわち、
\[
\mathcal{B}_{\mathrm{val}} \hookrightarrow \mathcal{A}
\]
という部分圏として、古典論理の二値性（真・偽）が、より広い論理的・計算的構造 $\mathcal{A}$ の中へ埋め込まれる。

\subsection{冗における役割}
冗において、$\mathcal{B}_{\mathrm{val}}$ は古典的な寓が扱う静的関係の層に対応する。したがって、冗は寓を否定するのではなく、それを二値截断として含み、その外側に動的評価層を拡張する。

この意味で、二値截断圏は「古典的数学が見ている部分」を明示する装置であり、視点数や継続残余は「古典的数学が截断した後に残る差分」を記述する装置である。

\subsection{まとめ}
二値截断圏 $\mathcal{B}_{\mathrm{val}}$ は、古典論理の二値性をより広い論理空間への埋め込みとして表現する部分圏であり、既存の数学的知識（部分圏・真理値・埋め込み）の枠内で一貫して説明可能である。

\section{中道不動点 \texorpdfstring{$\dfix$}{dfix}：\texorpdfstring{$E_\infty$}{E∞} における最大固定点}

\subsection{定義}
冗の公理系において、中道不動点 $\dfix$ は、「系全体の最大固定点である中道不動点 $\dfix$（空）へと収束する」と定義される。これは、動的プロセスが極限において収束する点を指す。

より正確には、$\dfix$ は通常の停止値ではない。むしろ、非停止性、自己言及、無限退行のように、二値評価では決定できない成分が、残余空間を経由して到達する極限的な固定点である。

\subsection{\texorpdfstring{$E_\infty$}{E∞} による説明}
視点数付き寓および冗の文脈において、$E_\infty$ は「極限同値（limit equivalence）」を表す記号として導入される。すなわち、動的軌道が有限段階では一致しないにもかかわらず、極限において同一視される構造を表す。

中道不動点 $\dfix$ は、この $E_\infty$ における最大固定点（greatest fixed point）として説明できる。具体的には、単調写像 $F$ に対して、
\[
\dfix = \bigvee \{ x \mid x \le F(x) \}
\]
という形で、動的プロセスが極限 $E_\infty$ において収束する点として定式化される。

\subsection{無記および空との関係}
冗では、二値截断圏で評価できない成分を、ただちに矛盾や失敗として扱わない。そのような成分は、無記として残余空間へ隔離される。

この隔離は、評価を停止させるためではなく、評価不能な成分を二値層から分離し、後続の動的過程を継続可能にするための操作である。極限において、この残余を伴う過程が到達する最大固定点が $\dfix$ であり、これは冗における「空」に対応する。

\subsection{まとめ}
中道不動点 $\dfix$ は、$E_\infty$ における最大固定点として説明可能であり、不動点理論の枠組みの中で一貫して扱うことができる。ただし冗における独自性は、この固定点を非停止プロセスの失敗ではなく、無記を経由した極限状態として読む点にある。

\section{本稿と冗の本論文の関係}
本稿は冗の公理系そのものを与えるものではなく、その前提となる概念を整理する準備論文である。したがって、本稿で扱う対象は、二値截断圏 $\mathcal{B}_{\mathrm{val}}$ と中道不動点 $\dfix$ に限定される。

冗の本論文では、これらを前提として、以下の構造を定式化する：
\begin{itemize}
\item 視点数 $p \in \mathbb{N} \cup \{\infty\}$ を伴う射 $R: X \to_{\langle p, \rho \rangle} Y$、
\item 交わり $R \cap S$ を観測文脈の動的併置として読む演算、
\item 逆関係 $R^\circ$ を未来制約の引き戻しとして読む双方向推論、
\item モジュラー律を継続残余の事前圧縮として読む残余モジュラー定理、
\item 非停止プロセスを無記として残余空間へ隔離し、$\dfix$ へ収束させる公理。
\end{itemize}

この分担により、本稿は既存数学との差分を明確にし、本論文は冗の新しい代数的・計算意味論的構造に集中できる。

\section{結論}
本稿では、
\begin{itemize}
\item 二値截断圏 $\mathcal{B}_{\mathrm{val}}$ を、古典論理の二値性がより広い論理空間へ埋め込まれる部分圏として説明し、
\item 中道不動点 $\dfix$ を、$E_\infty$ における最大固定点として説明した。
\end{itemize}
これにより、視点数付き寓および冗の概念構造が、一般的な数学知識の枠内で一貫して位置づけられることを示した。

この準備の上に、Freyd--Scedrov 型の寓に視点数と継続残余を加えた冗の公理系は定式化される。

\end{document}
